% Сборка из корня репозитория:  tectonic -X compile report/REPORT.tex --outdir .
\documentclass[a4paper,12pt]{article}

\usepackage{amsmath}
\usepackage{fontspec}
\usepackage{polyglossia}
\setdefaultlanguage{russian}
\setotherlanguage{english}
\setmainfont[Ligatures=TeX]{PT Serif}
\setsansfont[Ligatures=TeX]{PT Sans}
\setmonofont[Scale=0.9]{PT Mono}
\newfontfamily\cyrillicfont[Ligatures=TeX]{PT Serif}
\newfontfamily\cyrillicfontsf[Ligatures=TeX]{PT Sans}
\newfontfamily\cyrillicfonttt[Scale=0.9]{PT Mono}

\usepackage[left=30mm,right=15mm,top=20mm,bottom=20mm]{geometry}
\usepackage{setspace}
\onehalfspacing
\usepackage{indentfirst}
\setlength{\parindent}{1.25cm}
\usepackage{microtype}
\usepackage{float}
\setlength{\emergencystretch}{3em}
\sloppy
\usepackage{xcolor}
\definecolor{itmoblue}{HTML}{1946BA}
\usepackage{booktabs}
\usepackage{longtable}
\usepackage{tabularx}
\usepackage{array}
\usepackage{multirow}
\usepackage{enumitem}
\setlist{nosep,leftmargin=*}
\usepackage{titlesec}
\titleformat{\section}{\normalfont\large\bfseries}{\thesection.}{0.5em}{}
\titleformat{\subsection}{\normalfont\normalsize\bfseries}{\thesubsection.}{0.5em}{}
\titlespacing*{\section}{0pt}{1.2em}{0.6em}
\titlespacing*{\subsection}{0pt}{0.9em}{0.4em}
\usepackage{fancyhdr}
\pagestyle{fancy}
\fancyhf{}
\renewcommand{\headrulewidth}{0pt}
\fancyfoot[C]{\thepage}
\usepackage{fvextra}
\usepackage{caption}
\captionsetup{labelsep=endash,font=small}
\captionsetup[table]{justification=raggedright,singlelinecheck=false}
\usepackage[hidelinks]{hyperref}
\usepackage{url}
\usepackage{tikz}
\usetikzlibrary{arrows.meta,positioning}

\newcommand{\code}[1]{\texttt{#1}}
\newcolumntype{Y}{>{\raggedright\arraybackslash}X}
\newcolumntype{R}{>{\raggedleft\arraybackslash}X}

\begin{document}

% ---------------------------------------------------------------- Титульный лист
\begin{titlepage}
\begin{center}
\singlespacing
{\large Министерство науки и высшего образования Российской Федерации}\\[2mm]
{\large федеральное государственное автономное образовательное учреждение высшего образования}\\[2mm]
{\large\bfseries «Национальный исследовательский университет ИТМО»}\\[2mm]
{\color{itmoblue}\rule{\linewidth}{1pt}}

\vfill

{\Large\bfseries ОТЧЁТ ПО ПРОЕКТУ}\\[6mm]
{\LARGE\bfseries Ассистент сервисной службы\\[2mm] на Graph RAG}\\[6mm]
{\large Учебный кейс «ПолярТех»}

\vfill

\begin{flushright}
\begin{minipage}{0.6\linewidth}
\raggedright
\textbf{Выполнили:}\\[1mm]
Щетинкин Сергей Олегович\\
ИСУ 564310, \href{mailto:samara.serejka@gmail.com}{\code{samara.serejka@gmail.com}}\\[3mm]
Иванов Роман Евгеньевич\\
ИСУ 563976, \href{mailto:romanivanov0401@gmail.com}{\code{romanivanov0401@gmail.com}}\\[5mm]
\textbf{Репозиторий:}\\
{\small\href{https://github.com/byoverr/ai-thinking-project}{\code{github.com/byoverr/ai-thinking-project}}}
\end{minipage}
\end{flushright}

\vfill

Санкт-Петербург\\
2026
\end{center}
\end{titlepage}

\setcounter{page}{2}
\tableofcontents
\newpage

% ================================================================ Введение
\section*{Введение}
\addcontentsline{toc}{section}{Введение}

В этом проекте мы делали ассистента для сервисной службы учебной компании «ПолярТех», которая обслуживает холодильное оборудование. Ассистент должен отвечать на вопросы диспетчера по документации (руководствам, бюллетеням, регламентам и отчётам о выездах), узнавать цену и остаток детали в ERP и при необходимости оформлять заказ.

\textbf{Цель работы:} реализовать ассистента на основе Graph RAG и сравнить его с обычным (плоским) RAG на одном наборе вопросов.

Для этого нужно было решить следующие задачи:
\begin{enumerate}
\item реализовать извлечение сущностей и связей из документов с помощью LLM;
\item построить по ним граф знаний и разбить его на сообщества;
\item реализовать два режима поиска по графу: локальный и глобальный;
\item реализовать агента, который вызывает поиск и инструменты ERP и не создаёт заказ без подтверждения;
\item прогнать оценку всех режимов и разобрать, где граф помогает, а где нет.
\end{enumerate}

Инфраструктуру (клиент LLM, ERP, плоский RAG, оценщик, тесты платформы) мы получили готовой в шаблоне и не меняли. Наша часть кода находится в папках \code{src/assistant/graphrag/} и \code{src/assistant/agent/}.

% ================================================================ 1
\section{Постановка задачи}

Работа диспетчера обычно выглядит так: клиент сообщает модель оборудования и код ошибки, а диспетчеру нужно понять, какую деталь менять, есть ли она на складе, сколько стоит, и оформить заказ. Сложность в том, что нужная информация разбросана по разным документам. Например, в руководстве по модели указана исходная деталь, а то, что её сняли с производства и заменили на другую, написано только в сервисном бюллетене. Цену и остаток вообще знает только ERP.

Хороший пример вопроса, на котором ошибается плоский RAG: \emph{«Какую деталь заказать при ошибке E-17 на AR-210?»} Чтобы ответить правильно, нужно пройти по цепочке фактов из двух разных документов (рис.~\ref{fig:chain}).

\begin{figure}[H]
\centering
\begin{tikzpicture}[
  node distance=14mm and 26mm,
  ent/.style={draw, rounded corners, minimum height=7mm, inner sep=4pt, font=\small},
  rel/.style={-{Stealth}, thick},
  lab/.style={font=\scriptsize, fill=white, inner sep=1pt}]
\node[ent] (m) {AR-210};
\node[ent, right=of m] (p) {P-1042};
\node[ent, right=of p] (pb) {P-1042B};
\node[ent, below=of p] (c) {пусковое реле};
\node[ent, left=of c] (e) {E-17};
\draw[rel] (m) -- node[lab, above] {HAS\_PART} (p);
\draw[rel, red!70!black] (p) -- node[lab, above] {REPLACED\_BY} (pb);
\draw[rel] (p) -- node[lab, right] {IS\_A} (c);
\draw[rel] (e) -- node[lab, below] {INDICATES} (c);
\node[font=\scriptsize, gray, above=1mm of pb] {бюллетень SB-2025-03};
\node[font=\scriptsize, gray, above=1mm of m] {руководство AR-210};
\end{tikzpicture}
\caption{Цепочка фактов для вопроса про E-17 на AR-210. Красная связь есть только в бюллетене}
\label{fig:chain}
\end{figure}

Из руководства понятно, что E-17 указывает на пусковое реле, а реле в AR-210 это деталь P-1042. Но P-1042 снята с производства и заменена на P-1042B, и об этом написано только в бюллетене SB-2025-03. Поиск по сходству находит руководство (в нём есть и модель, и код ошибки), а бюллетень в выдачу не попадает. Поэтому плоский RAG советует заказать старую деталь. Граф знаний должен решать эту проблему: если факты из разных документов связаны через общие сущности, то цепочку можно пройти по рёбрам графа.

\paragraph{Данные.} Корпус состоит из 80 документов, после разбиения получилось 92 чанка. В схеме семь типов сущностей (линейка, модель, узел, деталь, код ошибки, регламент, поставщик) и восемь типов связей. Для проверки у нас было 24 открытых вопроса пяти типов. Скрытые вопросы, на которых проект будут проверять, нам недоступны, поэтому все числа в отчёте получены на открытых.

% ================================================================ 2
\section{Описание решения}

Общая схема системы показана на рис.~\ref{fig:pipeline}. Сначала документы один раз индексируются: LLM извлекает из каждого чанка сущности и связи, из них строится граф, граф делится на сообщества, и для каждого сообщества пишется сводка. Потом на вопрос можно отвечать в одном из трёх режимов: локальный поиск, глобальный поиск или агент.

\begin{figure}[H]
\centering
\begin{tikzpicture}[
  node distance=5mm and 7mm,
  box/.style={draw, rounded corners, align=center, font=\small, minimum height=9mm, inner sep=3pt},
  arr/.style={-{Stealth}, thick}]
\node[box] (d) {Документы\\(92 чанка)};
\node[box, right=of d] (x) {Извлечение\\(LLM)};
\node[box, right=of x] (g) {Граф\\(NetworkX)};
\node[box, right=of g] (s) {Сообщества\\и сводки};
\node[box, below=9mm of g] (l) {Local search};
\node[box, below=9mm of s] (gl) {Global search};
\node[box, below=of l, xshift=12mm] (a) {Агент + ERP};
\draw[arr] (d) -- (x);
\draw[arr] (x) -- (g);
\draw[arr] (g) -- (s);
\draw[arr] (g) -- (l);
\draw[arr] (s) -- (gl);
\draw[arr] (l) -- (a);
\draw[arr] (gl) -- (a);
\end{tikzpicture}
\caption{Схема системы}
\label{fig:pipeline}
\end{figure}

\subsection{Извлечение сущностей и связей}

Для каждого чанка мы отправляем модели запрос с JSON-схемой ответа, описанием направлений связей и несколькими примерами. Ответ проверяется через Pydantic. Если модель вернула невалидный JSON, делается ещё одна попытка, а если и она не удалась, чанк пропускается.

Маленькая модель довольно часто ошибается, поэтому мы не доверяем её ответу и фильтруем его:
\begin{itemize}
\item коды (модели, детали, ошибки) должны встречаться в тексте чанка после нормализации, иначе сущность отбрасывается. Нормализация нужна, потому что в отчётах коды пишут небрежно: \code{ar 210}, \code{AR210}, \code{Е-17} с русской «Е»;
\item названия узлов и линеек сверяются со словарём;
\item связь попадает в граф, только если типы её концов совпадают с сигнатурой из схемы (например, \code{HAS\_PART} может идти только от модели к детали).
\end{itemize}

\subsection{Граф и сообщества}

Граф мы строим в NetworkX как ориентированный. Узлы называются в виде \code{Тип:Имя}, а на узлах и рёбрах хранится список чанков, откуда взят факт, чтобы потом показывать источники ответа. Вес ребра равен числу \emph{разных} чанков, в которых встретился факт. Число повторов внутри одного ответа мы не учитываем, потому что модель иногда выдаёт одну и ту же связь несколько раз, и вес тогда был бы завышен.

Отдельно пришлось обработать встречные замены: иногда модель путает направление и выдаёт \code{P-1042B~REPLACED\_BY~P-1042}. Если в графе есть оба направления, мы оставляем то, у которого больше вес, а при равенстве удаляем оба, потому что непонятно, какое верное.

Для сообществ используется алгоритм Лувена (\code{seed=7}, чтобы результат воспроизводился). Перед этим граф переводится в неориентированный, и веса встречных рёбер складываются. Для каждого сообщества LLM пишет сводку до 150 слов. Числа в сводке (например, сколько было отказов узла) считает наш код и передаёт модели готовыми, так как сама модель считает плохо.

\subsection{Локальный и глобальный поиск}

\textbf{Локальный поиск} нужен для вопросов про конкретную модель, деталь или ошибку. Сначала в вопросе находятся коды и известные названия, это делается без LLM. Если таких опорных узлов несколько, мы берём все простые пути между ними длиной до трёх рёбер. Если опорный узел один, берём его окрестность глубины два. Дополнительно всегда добавляем цепочки \code{REPLACED\_BY} (с защитой от циклов), а если спрашивают про поставщика, то и связи \code{SUPPLIED\_BY}. Найденные факты передаются модели как контекст, его размер ограничен 8000 символами.

Мы заметили, что модель иногда находит в контексте замену, но в ответе всё равно пишет старую деталь. Поэтому в вопросах о детали, если конечная замена потерялась, мы восстанавливаем её из фактов графа. Также все коды в ответе проверяются: они должны быть либо в вопросе, либо в выбранных фактах.

\textbf{Глобальный поиск} нужен для обзорных вопросов вроде «какой узел чаще всего отказывает в линейке „Арктика“». Он работает по схеме map-reduce: каждую сводку сообщества модель оценивает отдельно и возвращает JSON \code{\{answer, score\}}, затем ответы с ненулевой оценкой объединяются одним reduce-запросом.

\subsection{Агент}

Агенту доступны пять инструментов: \code{search\_knowledge} (локальный поиск), \code{summarize\_topic} (глобальный поиск), \code{get\_price}, \code{check\_stock} и \code{create\_service\_order}. За один вопрос он может сделать не больше шести вызовов.

Основная идея в том, что все важные проверки делает код, а не модель:
\begin{itemize}
\item номер детали в аргументах должен встречаться в вопросе или в результатах предыдущих инструментов, иначе вызов отклоняется;
\item если для детали уже найдена замена, а модель пытается заказать старую, вызов тоже отклоняется;
\item одинаковые вызовы не повторяются, берётся прошлый результат;
\item если ERP недоступна, запрос повторяется (до трёх попыток);
\item в финальном ответе проверяются коды, цены и остатки. Если что-то не подтверждено, модель получает одну попытку исправить ответ, а после этого агент отказывается отвечать.
\end{itemize}

\subsection{Модель и запуск}

Все прогоны мы делали на \code{qwen3.5:4b} через Ollama (MLX, квантование nvfp4), эмбеддинги для плоского RAG считались моделью \code{bge-m3}. В README шаблона предлагалась другая модель (Qwen3~4B Q4\_K\_M), поэтому ориентир для baseline около 0,50 из задания к нашему прогону напрямую не относится.

Qwen~3.5 является thinking-моделью, и с включённым рассуждением она работала очень медленно: в пробном запросе на короткий JSON ушло 4768 выходных токенов. Чтобы это отключить, не меняя \code{LLMClient} из шаблона, мы написали небольшой HTTP-адаптер (\code{ollama\_gateway.py}). Он стоит между клиентом и Ollama и добавляет в каждый запрос \code{reasoning\_effort=none}, а ответы пропускает без изменений. Все итоговые прогоны сделаны без thinking.

% ================================================================ 3
\section{Результаты}

Все четыре режима мы оценивали стандартным оценщиком из шаблона на одних и тех же 24 вопросах. Ответ считается верным, если в нём есть все ключевые факты из эталона. Для вопросов с инструментами дополнительно проверяется, какие инструменты вызывались и было ли подтверждение. Стоит учитывать, что метрика строковая: она проверяет наличие нужных слов, а не смысл ответа.

\begin{table}[H]
\centering\small
\caption{Доля верных ответов по типам вопросов (в скобках число вопросов)}\label{tab:acc}
\setlength{\tabcolsep}{4pt}
\begin{tabular}{@{}lrrrrrr@{}}
\toprule
\textbf{Режим} & \textbf{Всего} & \textbf{Local (5)} & \textbf{Multi-hop (7)} & \textbf{Global (4)} & \textbf{Tool (5)} & \textbf{Refusal (3)}\\
\midrule
Плоский RAG  & 50,0\% & 100,0\% & 57,1\% & 0,0\%  & 0,0\%  & 100,0\%\\
Graph local  & 54,2\% & 100,0\% & 57,1\% & 25,0\% & 0,0\%  & 100,0\%\\
Graph global & 12,5\% & 20,0\%  & 0,0\%  & 0,0\%  & 0,0\%  & 66,7\%\\
Агент        & \textbf{62,5\%} & 80,0\% & 57,1\% & 0,0\% & 80,0\% & 100,0\%\\
\bottomrule
\end{tabular}
\end{table}

Из табл.~\ref{tab:acc} видно, что лучше всего показал себя агент (62,5\%), в основном за счёт вопросов с инструментами: остальные режимы ERP не вызывают и на такие вопросы ответить не могут. Локальный поиск по графу немного обошёл плоский RAG (54,2\% против 50,0\%), разница в одном вопросе. На multi-hop вопросах все режимы показали одинаковые 57,1\%, то есть граф здесь помог меньше, чем мы ожидали. Причины разбираются в разделах~\ref{sec:graph} и~\ref{sec:errors}. Глобальный поиск получился самым слабым: на конкретных вопросах сводки сообществ слишком общие, а на обзорных он не набрал ни одного верного ответа.

\begin{table}[H]
\centering\small
\caption{Качество источников и стоимость ответов}\label{tab:cost}
\begin{tabularx}{\linewidth}{@{}lRRRR@{}}
\toprule
\textbf{Режим} & \textbf{Recall источн.} & \textbf{Выдуманные факты} & \textbf{Вызовов LLM на вопрос} & \textbf{Токенов на вопрос}\\
\midrule
Плоский RAG  & 85,6\% & 0,0\% & 1,00 & 1115,0\\
Graph local  & 90,6\% & 0,0\% & 0,83 & 682,0\\
Graph global & 60,9\% & 0,0\% & 8,62 & 3174,6\\
Агент        & 90,6\% & 0,0\% & 5,75 & 4649,6\\
\bottomrule
\end{tabularx}
\end{table}

По стоимости (табл.~\ref{tab:cost}) локальный поиск оказался самым дешёвым: в среднем 0,83 вызова и 682 токена на вопрос. Вызовов меньше одного, потому что часть ответов (например, отказы, когда в графе ничего не найдено) формируется без обращения к модели. Агент дороже всех, около 4650 токенов на вопрос, так как он делает несколько шагов. Выдуманных фактов не нашлось ни в одном режиме, но эта метрика ловит только коды вне корпуса и цены или остатки, которые не пришли из ERP.

Токены считались как сумма входных и выходных по полю \code{usage} от сервера. Запросы из кеша оценщик считает как обычные, эмбеддинги в число вызовов не входят. Доля отклонённых вызовов агента составила 6,5\%, сюда же входят остановки перед созданием заказа без подтверждения.

\subsection{Извлечение и индексация}

\begin{table}[H]
\centering\small
\caption{Качество извлечения на эталоне из 16 документов}\label{tab:extr}
\begin{tabular}{@{}lrrrr@{}}
\toprule
\textbf{Объект} & \textbf{Precision} & \textbf{Recall} & \textbf{Найдено} & \textbf{В эталоне}\\
\midrule
Сущности & 100,0\% & 97,6\% & 163 & 165\\
Связи    & 100,0\% & 87,0\% & 142 & 161\\
\bottomrule
\end{tabular}
\end{table}

Precision получился 100\%, то есть всё, что прошло фильтр, совпало с эталоном. Recall связей ниже (87\%): часть верных связей модель не нашла, а часть была отброшена фильтром из-за перепутанного направления (подробнее в разделе~\ref{sec:errors}).

В итоговом графе 114 узлов, 255 рёбер и 8 сообществ. Индексация потребовала 100 вызовов модели: 85\,429 входных и 48\,024 выходных токенов. Конфликтов встречных замен на этом прогоне не возникло. Фильтр отбросил 18 сущностей и связей на всём корпусе и 16 на эталонных документах.

% ================================================================ 4
\section{Где граф помогает и где уступает}\label{sec:graph}

\subsection{Пример, где граф помогает (dev-06)}

Вопрос: \emph{«Какую деталь заказать при ошибке E-17 на AR-210?»} Это тот самый пример из раздела~1.

\begin{itemize}
\item \textbf{Плоский RAG (неверно):} «При ошибке E-17 на витрине AR-210 необходимо заменить деталь P-1042». В источниках только руководства \code{manual\_AR-210}, \code{manual\_AR-310} и \code{manual\_AR-220}, бюллетеня с заменой среди них нет.
\item \textbf{Graph local (верно):} «При ошибке E-17 на модели AR-210 необходимо заказать деталь P-1042B (замена для старой детали P-1042)».
\end{itemize}

Здесь граф сработал так, как и задумывалось. Связь \code{REPLACED\_BY} из бюллетеня и связи из руководства сходятся на узле P-1042, и локальный поиск прошёл всю цепочку с рис.~\ref{fig:chain}. Глобальный поиск на этом вопросе ошибся: в сводке сообщества отдельная деталь просто теряется.

\subsection{Пример, где граф уступает (dev-08)}

Вопрос: \emph{«Кто сейчас поставляет деталь узла „реле давления“ для модели GS-520?»}

\begin{itemize}
\item \textbf{Плоский RAG (верно):} деталь P-3126 снята с производства, вместо неё ставится P-3126B, поставщик «АвтоматикаСевер». В этот раз поиск по сходству нашёл бюллетень \code{SB-2025-07}.
\item \textbf{Graph local (неверно):} «Деталь для узла „реле давления“ в модели GS-520 это P-3126. Её поставщик: ХолодКомплект».
\end{itemize}

Мы посмотрели граф и выяснили, что при извлечении модель перепутала направление замены: в графе есть ребро \code{P-3126B~REPLACED\_BY~P-3126}, а правильного \code{P-3126~REPLACED\_BY~P-3126B} нет. Тип связи (деталь~$\to$~деталь) корректный, поэтому фильтр её пропустил, а встречного ребра, с которым можно было бы сравнить вес, тоже не было. В итоге поиск честно прошёл по графу, но сам граф оказался неверным. Получается, что Graph RAG очень зависит от качества извлечения: если факт извлечён с ошибкой, поиск это уже не исправит. Плоский RAG в этом случае работает с исходным текстом бюллетеня, поэтому ошибки извлечения его не затрагивают.

% ================================================================ 5
\section{Анализ ошибок}\label{sec:errors}

Ниже разобраны ошибки, которые мы нашли при просмотре ответов. Для каждой указан этап, на котором она возникла.

\begin{table}[H]
\centering\footnotesize
\caption{Разобранные ошибки системы}\label{tab:errors}
\begin{tabularx}{\linewidth}{@{}>{\raggedright\arraybackslash}p{1.3cm}>{\raggedright\arraybackslash}p{2.3cm}YY@{}}
\toprule
\textbf{Вопрос} & \textbf{Этап} & \textbf{Что произошло} & \textbf{Причина}\\
\midrule
dev-04 & Генерация (агент) & На вопрос «В какую линейку входит BZ-210?» агент ответил «в линейку моделей серии BZ», а в эталоне «Бриз». & Нужные документы найдены, но модель не назвала линейку так, как она называется в документах.\\
\addlinespace
dev-08 & Извлечение (агент) & Агент назвал заменой деталь P-3158 и не нашёл поставщика. & В графе нет связи \code{P-3126 REPLACED\_BY P-3126B} (см. раздел~\ref{sec:graph}).\\
\addlinespace
dev-09 & Извлечение (агент) & «Какая деталь ставится в BZ-110 как нагреватель оттайки?» Агент ответил, что номера детали в базе нет. Правильно: P-2054C. & В графе нет связи \code{BZ-110 HAS\_PART P-2054}.\\
\addlinespace
сообщ.~5 & Генерация сводки & В сводке написано, что у линейки «Арктика» 12 отказов, хотя по переданным весам их 14. & Числа считает код, но модель переписала их в тексте неточно, а отдельной проверки чисел в сводке у нас нет.\\
\bottomrule
\end{tabularx}
\end{table}

Чтобы понять, откуда берутся пропуски связей (dev-09), мы написали отдельный аудит (\code{audit.py}). Он заново прогоняет извлечение по кешу ответов модели и показывает, что именно отбросил фильтр. Всего было отброшено 17 связей и 1 сущность, ещё в 4 чанках модель вернула пустой результат. Самые показательные случаи:
\begin{itemize}
\item в \code{manual\_BZ-110} отброшено 8 связей, например \code{P-2010 HAS\_PART BZ-110}. Модель перевернула направление (деталь~$\to$~модель вместо модель~$\to$~деталь), и связь не прошла проверку схемы. Из-за этого и возникла ошибка dev-09;
\item в \code{manual\_GS-410} то же самое: 8 связей с перевёрнутым направлением;
\item в \code{report\_R-2025-024} модель склеила код модели с номером отчёта, такой код не нашёлся в тексте, и связь была отброшена.
\end{itemize}

Получается, что фильтр по схеме работает в обе стороны. Он не пускает в граф мусор (поэтому Precision 100\%), но вместе с мусором выбрасывает и правильные факты, у которых модель просто перепутала направление. Такие связи можно было бы не выбрасывать, а разворачивать, если после разворота они подходят под схему. Мы не успели это проверить.

\paragraph{Ограничения метрик.} Оценщик не проверяет каждое утверждение в ответе. Например, если агент назовёт деталь, которая есть в корпусе, но не подходит к вопросу, это не засчитается как выдуманный факт. Кроме того, вес \code{FAILED} считается по чанкам, из которых связь удалось извлечь, поэтому каждый пропущенный отчёт о выезде немного искажает статистику отказов.

% ================================================================ 6
\section{Границы системы}

Единственное действие, которое меняет данные в ERP, это создание заказа, и его ассистент без подтверждения человека не выполняет. Вот как это сделано:
\begin{itemize}
\item пока \code{confirm=False}, вызов \code{create\_service\_order} в ERP не уходит. Вместо этого агент возвращает \code{needs\_confirmation=True} и \code{pending\_action}, где указаны деталь, склад, количество и модель, то есть ровно то, что будет заказано;
\item даже при \code{confirm=True} количество должно быть явно названо в вопросе. Есть отдельная проверка, что цифры из кодов (например, 210 из AR-210) не принимаются за количество;
\item склад должен быть одним из четырёх (Москва, Санкт-Петербург, Казань, Новосибирск), отрицательное или дробное количество отклоняется, модель оборудования сверяется с вопросом;
\item на вопрос «какую деталь заказать?» заказ не создаётся, потому что это просьба о рекомендации. Заказ готовится только на явное «оформи» или «закажи»;
\item если ERP недоступна, агент так и пишет, а не выдаёт это за нулевой остаток.
\end{itemize}

Ещё мы столкнулись с тем, что на полностью заданный заказ модель иногда просто отвечает текстом «Подтверждаете?» и не вызывает инструмент. Тогда \code{pending\_action} не формируется. Мы добавили одно напоминание: если в вопросе есть всё для заказа, агент просит модель вызвать инструмент, и тот сам останавливается перед ERP.

Основные проверки покрыты тестами в \code{agent/tests/test\_safety.py}. Конечно, в учебном API подтверждение задаётся просто флагом \code{confirm}. В реальной системе подтверждение нужно было бы привязывать к конкретному \code{pending\_action}, а ERP нужен ключ идемпотентности, чтобы при сетевом сбое заказ не создался дважды.

% ================================================================ 7
\section{Следующий шаг}

Анализ ошибок показал, что основная проблема находится на этапе извлечения, а не поиска. Поэтому следующим шагом мы бы улучшали полноту извлечения, не меняя поиск. Идея такая: делать второй короткий проход только по тем чанкам, где в тексте есть коды, не попавшие в извлечённые сущности, или где для модели не нашлось связей \code{HAS\_PART} и \code{IS\_A}.

Проверить это можно на том же эталоне: сравнить Recall связей, точность на multi-hop вопросах и число токенов на индексацию. Будем считать идею удачной, если Recall и multi-hop вырастут, а Precision упадёт не больше чем на 0,02. После этого стоит проверить результат на новых документах, чтобы убедиться, что мы не подогнали систему под эталон.

% ================================================================ 8
\section{Вклад участников}

\begin{table}[H]
\centering\small
\begin{tabularx}{\linewidth}{@{}>{\raggedright\arraybackslash}p{3.4cm}Y@{}}
\toprule
\textbf{Участник} & \textbf{Что делал}\\
\midrule
Щетинкин\newline Сергей Олегович & Сообщества и сводки (\code{communities.py}), локальный и глобальный поиск (\code{search.py}), цикл агента (\code{agent.py}), адаптер для Ollama и тест к нему, тесты для graphrag (вместе с Романом), отчёт.\\
\addlinespace
Иванов\newline Роман Евгеньевич & Извлечение сущностей и связей (\code{extraction.py}), построение графа (\code{graph\_builder.py}), инструменты агента и проверки ответа (\code{tools.py}, \code{grounding.py}), тесты безопасности агента (вместе с Сергеем), индексация и итоговый прогон оценки, аудит извлечения.\\
\bottomrule
\end{tabularx}
\end{table}

% ================================================================ 9
\section{Использование ИИ-инструментов}

\begin{itemize}
\item \textbf{Codex} мы использовали, чтобы разобраться в задании и коде шаблона, а также как помощника при написании модулей и тестов, запуске экспериментов и разборе результатов.
\item \textbf{Claude Code} помогал с вёрсткой отчёта в \LaTeX.
\item \textbf{Qwen~3.5~4B} это рабочая модель самой системы: она извлекает сущности, пишет сводки, отвечает на вопросы и выбирает инструменты. \textbf{bge-m3} считает эмбеддинги для плоского RAG.
\end{itemize}

Все числа в отчёте взяты из файлов оценщика в папке \code{artifacts/}. Заглушки \code{FakeLLM} использовались только в тестах.

% ================================================================ Заключение
\section*{Заключение}
\addcontentsline{toc}{section}{Заключение}

В ходе работы мы реализовали ассистента сервисной службы на Graph RAG: извлечение сущностей и связей, граф знаний, сообщества, локальный и глобальный поиск и агента с инструментами ERP. Все 168 тестов проходят.

По результатам можно сделать следующие выводы:
\begin{enumerate}
\item Агент дал лучший результат (62,5\% верных ответов против 50,0\% у плоского RAG), в первую очередь потому, что только он умеет обращаться к ERP.
\item Локальный поиск по графу оказался немного точнее плоского RAG и при этом дешевле него примерно в 1,6 раза по токенам.
\item Граф хорошо помогает, когда ответ собирается из нескольких документов (пример dev-06), но качество ответов сильно зависит от извлечения. Одна перевёрнутая связь, как в dev-08, приводит к неверному ответу.
\item Глобальный поиск на нашем наборе вопросов работает плохо и стоит дороже остальных режимов, кроме агента.
\item Создание заказа надёжно защищено проверками в коде: без подтверждения и явно названного количества заказ не создаётся.
\end{enumerate}

Главное, что стоит улучшать дальше, это полноту извлечения связей, так как большинство разобранных ошибок возникает именно там.

% ================================================================ Приложение
\appendix
\section*{Приложение. Как воспроизвести результаты}
\addcontentsline{toc}{section}{Приложение. Как воспроизвести результаты}

Прогон делался на Linux ARM64 с Python 3.12 без Docker, по инструкции из \code{docs/setup.md}. В отдельных терминалах из корня проекта:

{\small
\begin{Verbatim}[breaklines]
uv sync --python 3.12
uv run uvicorn erp.app:app --port 8100
OLLAMA_UPSTREAM=http://<адрес-ollama>:11434/v1 uv run python -m assistant.agent.ollama_gateway
\end{Verbatim}
}

\noindent В файле \code{.env}:
{\small
\begin{Verbatim}
LLM_BASE_URL=http://127.0.0.1:11435/v1
LLM_MODEL=qwen3.5:4b
EMBED_MODEL=bge-m3
ERP_FAIL_RATE=0
\end{Verbatim}
}

\noindent Дальше:
{\small
\begin{Verbatim}
make index RUN="uv run"
make test RUN="uv run"
make lint RUN="uv run"
make eval-all RUN="uv run"
make eval-extraction RUN="uv run"
\end{Verbatim}
}

Кеш ответов модели лежит в \code{artifacts/llm\_cache/}, с ним прогон повторяется за несколько минут. Он относится к режиму без thinking. Кеш пробного запроса с thinking лежит отдельно в \code{artifacts/pilot\_thinking\_cache/} и в результатах не используется. Если менять режим thinking, кеш нужно держать в отдельной папке, потому что \code{LLMClient} не учитывает этот параметр в ключе кеша.


\noindent Отчёт собирается из \code{report/REPORT.tex} командой:
{\small
\begin{Verbatim}
tectonic -X compile report/REPORT.tex --outdir .
\end{Verbatim}
}

\end{document}
